\section*{Розділ VI. Прикінцеві положення}
\addcontentsline{toc}{section}{Розділ VI. Прикінцеві положення}

\subsection*{6.1. Затвердження та набуття чинності}
\addcontentsline{toc}{subsection}{6.1. Затвердження та набуття чинності}
    6.1.1. Цей Тимчасовий виборчий регламент затверджується Загальними зборами студентів Університету простою більшістю голосів від присутніх студентів за умови досягнення кворуму не менше 200 студентів або 15\% від загальної кількості студентів Університету (береться менша цифра).

    6.1.2. Цей Регламент набуває чинності з моменту його затвердження Загальними зборами студентів та офіційного оприлюднення на інформаційних ресурсах Університету.

    6.1.3. З моменту набуття чинності цим Регламентом припиняють дію всі попередні нормативні акти з питань проведення виборів до органів студентського самоврядування, що суперечать положенням цього Регламенту.

\subsection*{6.2. Тлумачення норм Регламенту}
\addcontentsline{toc}{subsection}{6.2. Тлумачення норм Регламенту}
    6.2.1. Офіційне тлумачення норм цього Регламенту у спірних випадках надає Тимчасова виборча комісія шляхом прийняття відповідних роз'яснень.

    6.2.2. Роз'яснення ТВК щодо застосування норм цього Регламенту є обов'язковими для всіх учасників виборчого процесу та можуть бути оскаржені до Загальних зборів студентів згідно з процедурою, встановленою Розділом V.

    6.2.3. У випадках, не передбачених цим Регламентом, ТВК керується загальними принципами демократичності, справедливості, прозорості та забезпечення виборчих прав студентів.

\subsection*{6.3. Припинення дії Регламенту}
\addcontentsline{toc}{subsection}{6.3. Припинення дії Регламенту}
    6.3.1. Дія цього Регламенту припиняється після виконання умови, зазначеної у пункті 1.4.2 Розділу I: проведення перших виборів делегатів КСУ та скликання першого засідання КСУ.

    6.3.2. Незалежно від виконання умов пункту 6.3.1, дія цього Регламенту автоматично припиняється через 6 місяців з дня набуття ним чинності.

    6.3.3. Про припинення дії цього Регламенту ТВК інформує студентську спільноту шляхом офіційного повідомлення не пізніше ніж за 3 (три) дні до дати припинення його дії.

    6.3.4. Після припинення дії цього Регламенту всі подальші вибори до органів студентського самоврядування проводяться за процедурами, встановленими Положенням про Центральну виборчу комісію студентів та іншими постійно діючими нормативними актами студентського самоврядування.

\subsection*{6.4. Перехідні положення}
\addcontentsline{toc}{subsection}{6.4. Перехідні положення}
    6.4.1. Формування Тимчасової виборчої комісії здійснюється на тих самих Загальних зборах студентів, що затверджують цей Регламент, згідно з процедурою, визначеною Розділом II.

    6.4.2. ТВК приступає до виконання своїх функцій протягом 3 днів після свого формування і проводить перше засідання для обрання керівництва та затвердження плану роботи.

    6.4.3. Перші вибори до органів студентського самоврядування повинні розпочатися не пізніше ніж через 2 тижні після формування ТВК.

    6.4.4. Адміністрація Університету зобов'язана надати ТВК необхідну організаційну, технічну та інформаційну підтримку для проведення перших виборів, включаючи:

        \begin{enumerate}[label=\alph*)]
            \item Надання приміщень для засідань ТВК та проведення голосування;
            \item Технічне забезпечення для організації електронного голосування (за потреби);
            \item Доступ до офіційних інформаційних ресурсів Університету;
            \item Актуальні списки студентів для формування виборчих округів;
            \item Інше організаційно-технічне забезпечення виборчого процесу.
        \end{enumerate}

    6.4.5. Всі витрати, пов'язані з проведенням перших виборів, покриваються за рахунок коштів Університету в межах видатків на забезпечення діяльності органів студентського самоврядування.

\subsection*{6.5. Документообіг та архівування}
\addcontentsline{toc}{subsection}{6.5. Документообіг та архівування}
    6.5.1. Усі документи, створені під час дії цього Регламенту (протоколи засідань ТВК, рішення, документи кандидатів, протоколи голосування тощо), ведуться в електронному форматі з обов'язковим створенням резервних копій.

    6.5.2. Документи підписуються кваліфікованими електронними підписами уповноважених осіб або, за неможливості, паперовими підписами з подальшим скануванням.

    6.5.3. Після припинення дії цього Регламенту вся документація та матеріали передаються:

        \begin{enumerate}[label=\alph*)]
            \item Центральній виборчій комісії студентів -- для використання в подальшій роботі;
            \item Студентській уповноваженій делегації -- для архівного зберігання як частина історії формування ОСС;
            \item Архіву Університету -- для постійного зберігання.
        \end{enumerate}

    6.5.4. Строк зберігання документів перших виборів становить не менше 5 років з моменту припинення дії цього Регламенту.

    6.5.5. Обробка персональних даних та технічних журналів, пов'язаних з проведенням перших виборів, здійснюється відповідно до вимог законодавства України про захист персональних даних.

\subsection*{6.6. Заключні положення}
\addcontentsline{toc}{subsection}{6.6. Заключні положення}
    6.6.1. Цей Тимчасовий виборчий регламент є тимчасовим нормативним актом студентського самоврядування, призначеним для одноразового використання з метою запуску системи органів студентського самоврядування Університету.

    6.6.2. Досвід застосування цього Регламенту може бути використаний при розробці постійних процедур організації виборів до ОСС та вдосконаленні нормативної бази студентського самоврядування.

    6.6.3. Цей Регламент сприяє розвитку демократичної культури в студентському середовищі та забезпечує формування легітимної та ефективної системи студентського самоврядування в Університеті.

    6.6.4. Питання, не врегульовані цим Регламентом, вирішуються ТВК з урахуванням принципів студентського самоврядування, демократичності та забезпечення прав студентів, визначених Положенням про органи студентського самоврядування Університету.